% ------------------------------------------------------------------------------
% Capítulo 2 · Antecedentes
% Fuentes: propuesta del TFG, docs/tfg/DECISIONES.md (ADR-001, 004, 005, 011,
% 014, 017, 020), bibliografía. Por ADR-007 no se nombran los componentes base
% utilizados.
% ------------------------------------------------------------------------------

Este capítulo describe la situación de la que parte el trabajo: qué necesitan
las pymes cuando quieren ofrecer o usar software SaaS propio, qué piezas
comparten casi todas esas aplicaciones y qué soluciones existen hoy para no
tener que construirlas desde cero. A partir de ese análisis se justifica la
solución adoptada.

\section{Contexto}
\label{sec:ctx-contexto}

\subsection{Las pymes y el software como servicio}
\label{sec:ctx-pymes-saas}

La Comisión Europea define como pyme a la empresa que emplea a menos de 250
personas y cuyo volumen de negocio anual no supera los 50 millones de euros o
cuyo balance general no supera los 43 millones~\cite{ec-sme-definition}. Son
la gran mayoría de las empresas, pero sus recursos técnicos son limitados:
pocas tienen un departamento de informática, y las que encargan una aplicación
a medida lo hacen con presupuestos y plazos ajustados.

El modelo SaaS encaja con esa situación. Según la definición del
NIST~\cite{nist-cloud}, en este modelo el proveedor ofrece una aplicación que
se ejecuta en una infraestructura en la nube gestionada por él; el cliente la
usa a través del navegador sin preocuparse de servidores, actualizaciones ni
copias de seguridad, y paga por suscripción. Para una pyme, el SaaS aparece en
dos papeles:

\begin{itemize}
  \item como \textbf{consumidora}, cuando contrata una herramienta en lugar de
        instalarla; y
  \item como \textbf{proveedora}, cuando su propio negocio se ofrece en
        forma de aplicación web: una zona privada para sus clientes, un
        servicio de reservas, un catálogo con suscripción o, simplemente, la
        web de la empresa con un acceso para clientes y otro para el
        personal.
\end{itemize}

Este trabajo se sitúa en el segundo caso, y en particular en el de quien
\emph{desarrolla} esas aplicaciones para pymes: un estudio pequeño o un
desarrollador independiente que debe entregar, con poco presupuesto, una
aplicación segura y mantenible, y que probablemente tendrá que repetir el
ejercicio para otro cliente poco después.

\subsection{Componentes comunes de una aplicación SaaS}
\label{sec:ctx-componentes}

Por diferente que sea el negocio, una aplicación SaaS repite un conjunto de
componentes transversales. La tabla~\ref{tab:ctx-componentes} los resume; son,
precisamente, los que la propuesta del trabajo enumera como esenciales.

\begin{table}[htbp]
  \centering
  \small
  \caption{Componentes que se repiten en casi cualquier aplicación SaaS.}
  \label{tab:ctx-componentes}
  \begin{tabularx}{\textwidth}{@{}>{\raggedright\arraybackslash}p{3.6cm}X@{}}
    \toprule
    \textbf{Componente} & \textbf{Qué debe resolver} \\
    \midrule
    Interfaz pública y privada & Página de presentación, precios y contenido; área privada con navegación y ajustes. \\
    Autenticación & Registro, verificación del correo, inicio de sesión con contraseña o enlace, recuperación de la contraseña, segundo factor (MFA) y, opcionalmente, proveedores externos. \\
    Cuentas y organizaciones & Cuenta personal, cuentas de organización con miembros, invitaciones, roles y permisos por funcionalidad. \\
    Aislamiento entre clientes & Garantizar que los datos de una organización no son visibles ni modificables desde otra (\textit{multi-tenancy}). \\
    Pagos & Planes, alta de la suscripción, portal del cliente y sincronización del estado con la pasarela de pago. \\
    Administración & Consola para el personal: buscar usuarios, bloquearlos, darles soporte, consultar métricas. \\
    Gestión de datos y contenidos & Consultar y corregir datos, publicar contenidos y gestionar archivos sin programar, con permisos y trazabilidad. \\
    Transversales & Correos transaccionales, notificaciones, internacionalización, registro de errores, pruebas y despliegue. \\
    \bottomrule
  \end{tabularx}
\end{table}

Ninguno de estos componentes diferencia un producto de otro, pero su coste
de desarrollo es alto y su margen de error, muy pequeño. La literatura sobre
aplicaciones multicliente insiste en que el aislamiento, la configurabilidad
por cliente y la seguridad son las preocupaciones de arquitectura que más
condicionan estos sistemas y las que más complican su
mantenimiento~\cite{bezemer2010multitenant,krebs2012multitenancy}. La
ingeniería del software lleva décadas proponiendo la misma respuesta a este
tipo de problema: identificar lo común a una familia de productos y
construirlo una vez, como una base reutilizable sobre la que se derivan los
productos concretos~\cite{clements2001spl}.

\section{Estado de la técnica}
\label{sec:ctx-estado-tecnica}

Se revisan a continuación las cuatro familias de soluciones que cubren, total
o parcialmente, los componentes anteriores: las plantillas o kits de
aplicaciones SaaS, los servicios de \textit{backend} gestionado, los gestores
de contenidos y paneles de administración de bases de datos, y las técnicas de
aislamiento multicliente en la base de datos.

\subsection{Plantillas y kits de aplicaciones SaaS}
\label{sec:ctx-kits}

Desde hace años existe un mercado de plantillas de arranque
(\textit{boilerplates} o \textit{starter kits}) para aplicaciones SaaS. Son
proyectos que ya incluyen la autenticación, las cuentas de equipo, la
facturación y una página de presentación, construidos sobre un
\textit{framework} web concreto, y que el desarrollador copia y adapta. Se
distribuyen tanto con licencias libres como con licencias comerciales de pago
único, y estas últimas suelen ser más completas: incorporan cuentas de equipo
con roles, panel de administración, pruebas automatizadas y documentación.

Su principal ventaja es el ahorro de tiempo inicial. Sus limitaciones, desde
el punto de vista de este trabajo, son tres:

\begin{enumerate}
  \item rara vez incluyen una herramienta para que el personal no técnico
        administre los datos, más allá de un panel de super-administración
        con unas pocas operaciones sobre usuarios;
  \item están pensadas para el mercado anglosajón: interfaz, correos,
        comentarios y documentación en inglés; y
  \item su calidad de seguridad no está garantizada. Al ser código ajeno,
        quien lo adopta asume también sus fallos, y la revisión corre de su
        cuenta.
\end{enumerate}

\subsection{Servicios de \textit{backend} gestionado (BaaS)}
\label{sec:ctx-baas}

Los servicios de \textit{Backend as a Service} ofrecen ya construidos la base
de datos, la autenticación, el almacenamiento de archivos y una API de acceso a
los datos. La tabla~\ref{tab:ctx-baas} compara los más extendidos en sus
características relevantes para este trabajo.

\begin{table}[htbp]
  \centering
  \small
  \caption{Comparación de servicios de \textit{backend} gestionado.}
  \label{tab:ctx-baas}
  \begin{tabularx}{\textwidth}{@{}l>{\raggedright\arraybackslash}X>{\raggedright\arraybackslash}X>{\raggedright\arraybackslash}p{2.6cm}@{}}
    \toprule
    \textbf{Servicio} & \textbf{Modelo de datos} & \textbf{Autorización} & \textbf{Código / alojamiento} \\
    \midrule
    Firebase~\cite{firebase-docs} & Documental (Cloud Firestore) & Reglas de seguridad propias, declaradas aparte de los datos & Propietario; solo en la nube del proveedor \\
    Supabase~\cite{supabase-docs} & Relacional (PostgreSQL) con API generada & Políticas RLS de PostgreSQL sobre la identidad del JWT~\cite{supabase-rls} & Abierto; nube o autoalojado \\
    Appwrite~\cite{appwrite-docs} & Colecciones de documentos sobre una base relacional interna & Permisos por colección y documento & Abierto; nube o autoalojado \\
    PocketBase~\cite{pocketbase} & SQLite embebido & Reglas de API por colección & Abierto; un único ejecutable autoalojado \\
    \bottomrule
  \end{tabularx}
\end{table}

Para una aplicación multicliente con relaciones entre cuentas, miembros,
roles, suscripciones y datos de negocio, un modelo relacional con integridad
referencial es una ventaja clara. Supabase añade otra: su autorización no es
un sistema paralelo, sino la propia seguridad a nivel de fila de
PostgreSQL~\cite{postgresql-rls}. Cada petición llega a la base de datos con
la identidad del usuario, y las políticas deciden qué filas puede ver o
modificar. Al ser PostgreSQL estándar, además, el esquema, las políticas y sus
pruebas pueden versionarse como código y probarse con herramientas como
pgTAP~\cite{pgtap}. Su contrapartida es que la seguridad depende de que esas
políticas estén bien escritas, y un error en ellas no se ve en la interfaz.

\subsection{Gestores de contenidos y paneles de administración de datos}
\label{sec:ctx-cms}

El componente de administración de datos admite varias soluciones:

\begin{itemize}
  \item \textbf{CMS \textit{headless}} como Strapi~\cite{strapi-docs}. El
        desarrollador define tipos de contenido y el CMS genera una API y un
        panel para editarlos. Son excelentes para contenidos editoriales,
        pero gestionan su propio esquema y su propio sistema de usuarios: no
        están pensados para administrar las tablas de negocio de otra
        aplicación ni para reutilizar sus políticas de acceso.
  \item \textbf{Plataformas de datos sobre una base existente} como
        Directus~\cite{directus-docs}, que se conectan a una base de datos SQL
        ya creada y ofrecen sobre ella una API y una interfaz de edición con
        roles y permisos. Se acercan más a lo que se busca, pero añaden un
        segundo servicio de \textit{backend} con su propio modelo de usuarios
        y permisos que hay que sincronizar con el de la aplicación.
  \item \textbf{Herramientas para desarrolladores}, como el editor de tablas
        del panel de Supabase~\cite{supabase-docs} o pgAdmin~\cite{pgadmin}.
        Son potentes, pero trabajan con privilegios de administrador de la
        base de datos, sin permisos por tabla para el personal y sin
        auditoría orientada al negocio. No se pueden entregar a quien gestiona
        una pyme.
  \item \textbf{CMS específicos para un BaaS}, que se instalan junto a un
        proyecto de Supabase y ofrecen un explorador de datos, de usuarios y de
        almacenamiento con un RBAC propio guardado en la misma base de datos.
        Son la opción más cercana al objetivo, aunque normalmente se
        distribuyen como una aplicación independiente, con su propio inicio de
        sesión y su propia pila tecnológica.
\end{itemize}

\subsection{Aislamiento multicliente en la base de datos}
\label{sec:ctx-multitenancy}

Existen tres estrategias clásicas para aislar los datos de los clientes de una
aplicación SaaS~\cite{chong2006multitenant}: una base de datos por cliente; un
esquema por cliente dentro de una base compartida; y tablas compartidas en las
que cada fila lleva el identificador del cliente al que pertenece. La primera
aísla mejor, pero multiplica el coste de operación; la última es la más
económica y la que mejor escala en número de clientes, a cambio de que el
aislamiento dependa por completo de que \emph{todas} las consultas filtren por
cliente~\cite{krebs2012multitenancy}.

La seguridad a nivel de fila resuelve esa debilidad trasladando el filtro de
la aplicación a la base de datos~\cite{postgresql-rls,postgresql-createpolicy}.
Cada tabla declara políticas del tipo «un usuario solo puede leer las filas de
las cuentas de las que es miembro», y PostgreSQL las aplica a cualquier
consulta, venga de donde venga. Un fallo en el código de la aplicación ya no
basta para filtrar datos de otro cliente: tendría que fallar también la
política. Es el enfoque que recomienda Supabase~\cite{supabase-rls} y el que
adopta este trabajo, complementado con pruebas automáticas que atacan las
políticas suplantando a usuarios de distintos clientes.

\section{Justificación de la solución}
\label{sec:ctx-justificacion}

Del análisis anterior se desprenden las decisiones que dan forma a PymeKit:

\begin{enumerate}
  \item \textbf{Una base reutilizable, no una aplicación.} Lo que se repite
        de un SaaS a otro es la infraestructura de la
        tabla~\ref{tab:ctx-componentes}. Construirla una vez, bien probada y
        configurable, es lo que acelera los proyectos siguientes (RNF-01).
  \item \textbf{PostgreSQL con RLS como frontera de seguridad.} Un modelo
        relacional encaja con los datos de un SaaS multicliente, y la RLS hace
        que el aislamiento entre clientes no dependa solo del código. Se
        adopta Supabase porque ofrece autenticación, almacenamiento y
        PostgreSQL estándar, es de código abierto y permite un entorno local
        completo para el desarrollo y las pruebas.
  \item \textbf{Un CMS de datos dentro de la propia aplicación.} Ni los CMS
        \textit{headless} ni las herramientas para desarrolladores cubren la
        necesidad de que el personal de una pyme administre los datos de
        \emph{su} aplicación con permisos y auditoría. Un CMS que trabaja
        sobre la misma base de datos sí lo hace, y si además se integra en la
        consola de administración se evita un segundo servicio, un segundo
        inicio de sesión y una segunda pila tecnológica (ADR-011).
  \item \textbf{Español e internacionalización.} Una base dirigida a pymes
        españolas debe estar en español, pero sin cerrar la puerta a otros
        idiomas (ADR-021).
\end{enumerate}

\paragraph{Partir de componentes base.}
Como decisión de ingeniería, PymeKit se construye a partir de componentes base
con licencia comercial ---una base de aplicación SaaS y un CMS para la base de
datos---, una práctica habitual en la industria y coherente con el propio
objetivo de reutilización de este trabajo (ADR-001). La aportación del trabajo
está en la arquitectura, la integración en una sola aplicación, la adaptación
a pymes españolas, la verificación de su seguridad y su reutilizabilidad. Que una base reutilizable deba verificarse
antes de reutilizarse es, en sí mismo, uno de los resultados que se analizan en
el capítulo de conclusiones.
